\documentclass[10pt,a4paper]{amsart}
\usepackage{amsmath,amssymb,amsthm,graphicx}
\usepackage{amsthm,amsmath, amsfonts, amssymb, comment}
\newtheorem{theorem}{Theorem}[section]
\newtheorem{lemma}[theorem]{Lemma}
\newtheorem{proposition}[theorem]{Proposition}
\newtheorem{corollary}[theorem]{Corollary}
\newtheorem{question}[theorem]{Question}
\newtheorem*{conjecture}{Conjecture}
\theoremstyle{definition}
\newtheorem{definition}{Definition}[section]
\newtheorem*{programme}{Program}
\theoremstyle{remark}
\newtheorem*{remark}{Remark}
\usepackage{enumerate}
\title{Triangle equivalences of distance value sets accumulating at zero are similarities}
\author{Mahmood Etedadialiabadi, Su Gao, Feng Li, Ruiwen Li}
\date{arxiv:2411.02760v2; CC BY 4.0}
\begin{document}
\maketitle
\noindent\textit{Editorial note.} The entire author proof is retained. At1028 the triangle estimate supplies <=f(delta\_i)<epsilon, sufficient for continuity. At1043 an attained finite upper endpoint is covered by the already established continuity from below; the independent review explains this boundary instance of the author limit argument.\par\medskip
\noindent\textit{Editorial scope.} One original rigidity theorem, including its assertion about every equivalence witness. Ordinary real analysis is allowed. Neither the author discrete-distance rigidity theorem nor the Borel order-coding theorem is a free premise or an additional counted item.\par\medskip
\noindent\textit{Allowed standard prerequisites (editorial).}\begin{itemize}
\item Elementary real analysis. Uniform continuity on intervals, density, rational approximation and limits.. Convert triangle constraints into continuity and then constant ratios; none of the rigidity conclusions is assumed.
\end{itemize}\medskip
    \begin{definition}[\cite{EGLMM}]
        Let $\Delta$ be a nonempty set of positive real numbers. We call $\Delta$ a {\em distance value set} if for all $x, y\in \Delta$, 
        \begin{equation*}
           \min\{ x+y, \sup\Delta\} \in \Delta.
        \end{equation*}
    \end{definition}
Briefly speaking, a distance value set is exactly an initial segment of a subsemigroup of positive reals. Given a distance value set $\Delta$, a metric space $(X,d)$ is a {\em $\Delta$-metric space} if $d(x,y)\in\Delta$ for all $x\neq y\in X$. 

\begin{definition}[\cite{EGLMM}]\label{def:CDVSequiv} Let $\Delta$ and $\Lambda$ be distance value sets. 
\begin{enumerate}
\item[(i)] A triple $(x,y,z)$ of positive real numbers is a {\em $\Delta$-triangle} if $x,y, z\in \Delta$ and $|x-y|\leq z\leq x+y$.
\item[(ii)] We say that $\Delta$ and $\Lambda$ are {\em equivalent} if there is a bijection $f\colon \Delta\to \Lambda$ such that for any positive real numbers $x,y, z$, $(x,y,z)$ is a $\Delta$-triangle if and only if $(f(x), f(y), f(z))$ is a $\Lambda$-triangle.
\end{enumerate}
\end{definition}

\begin{lemma}\label{lem:isomorphism} $\approx_{\mathsf{CDVS}}$ is Borel reducible to an isomorphism relation of countable structures. Consequently, $\approx_{\mathsf{CDVS}}$ is Borel reducible to $E^\infty_{S_\infty}$.
\end{lemma}

\begin{proof} To each $\Delta\in \mathsf{CDVS}$ we associate a countable structure $\mathcal{A}_\Delta=(\Delta, R^{\mathcal{A}})$, where $R^{\mathcal{A}}$ is a ternary relation defined as
$$ R^{\mathcal{A}}(d_1,d_2,d_3)\iff \mbox{$(d_1,d_2,d_3)$ is a $\Delta$-triangle.}$$
It is clear that for $\Delta, \Lambda\in \mathsf{CDVS}$, we have $\Delta\approx \Lambda$ if and only if $\mathcal{A}_\Delta\cong \mathcal{A}_{\Lambda}$. This gives a Borel reduction from $\approx_{\mathsf{CDVS}}$ to an isomorphism relation of countable structures.
\end{proof}

 \begin{theorem}\label{thm:simapprox}
        Let $\Delta$ and $\Lambda$ be countable distance value sets with $\inf\Delta=\inf\Lambda=0$. Then the following are equivalent: 
        \begin{enumerate}
            \item[\rm (1)] $\Delta$ and $\Lambda$ are equivalent. 
            \item[\rm (2)] There is a positive real number $r$ such that $\Lambda = r \cdot \Delta$. 
        \end{enumerate} 

        Moreover, if $\Delta$ and $\Lambda$ are equivalent, then every witness of their equivalence has the form of multiplication by a positive real number. 
     \end{theorem}

\begin{proof} The implication (2)$\Rightarrow$(1) is obvious. We only prove (1)$\Rightarrow$(2). For this, let $f$ be a bijection between $\Delta$ and $\Lambda$ witnessing their equivalence. Let $\mathcal{A}_\Delta=(\Delta, R^\Delta)$ and $\mathcal{A}_\Lambda$ be the countable structures defined in the proof of Lemma~\ref{lem:isomorphism}. Then $f$ is an isomorphism witnessing $\mathcal{A}_\Delta\cong \mathcal{A}_\Lambda$. Since $\Delta$ is a distance value set and $\inf\Delta=0$, we have that $\Delta$ is dense in $(0, \sup\Delta)$. Similarly, $\Lambda$ is dense in $(0, \sup\Lambda)$.

Fix a sequence $(\delta_i)_{i\in\mathbb{N}}$ in $\Delta$ such that $\delta_{i+1} < {\delta_i}/{2}$ for all $i$. Such a sequence exists since $\mathcal{A}_\Delta$ satisfies the sentence $\forall x\,\exists y\,\neg R(y,y,x)$. Since $f$ is an isomorphism, we have $f(\delta_{i+1})<{f(\delta_i)}/{2}$ for all $i$. 

We claim that $f$ is continuous as a function from $\Delta\subseteq\mathbb{R}$ to $\Lambda\subseteq\mathbb{R}$. To see this,  
fix $x \in \Delta$. For any $\epsilon > 0$, find an $i$ such that $\delta_i < x$ and $f(\delta_i) < \epsilon$. For any $y \in \Delta$ with $|x - y| < \delta_i$, $(x,y,\delta_i)$ is a $\Delta$-triangle. Thus $(f(x), f(y),f(\delta_i))$ is a $\Lambda$-triangle. In particular, $|f(x)-f(y)|<f(\delta_i)<\epsilon$. The actually proves that for all positive $\delta$, $f$ is uniformly continuous on $\Delta_{>\delta}\triangleq \Delta\cap (\delta,+\infty)$. 

Next we show that $f$ is in fact linear. In the proof we use notation such as $\Delta_{>\delta}$, $\Delta_{<\delta}$, and $\Delta_{[a,b]}$ with obvious meanings.

     
     %  If $x,y, x+y \in\Delta$, then $f(x+y) = f(x)+f(y)$. 
   
To show that $f$ is linear, assume $x,y,x+y\in \Delta$. We claim that
\begin{enumerate}
\item[(a)] $f(\Delta_{[|x-y|, x+y]}) \subseteq \Lambda_{[|f(x)-f(y)|, f(x)+f(y)]}$; 
\item[(b)] $f(\Delta_{>x+y}) \subseteq \Lambda_{>f(x)+f(y)}  \mbox{ or } f(\Delta_{>x+y})\subseteq \Lambda_{<|f(x)-f(y)|}$;
\item[(c)] $f(\Delta_{<|x-y|}) \subseteq \Lambda_{>f(x)+f(y)} \mbox{ or } f(\Delta_{<|x-y|})\subseteq \Lambda_{<|f(x)-f(y)|}$. 
\end{enumerate}
Let $z\in \Delta$. (a) is equivalent to saying that if $(x,y,z)$ is a $\Delta$-triangle, then $(f(x), f(y), f(z))$ is a $\Lambda$-triangle, which is immediate from our assumption. Also immediate is that if $z > x+y$, then $(x,y,z)$ is not a $\Delta$-triangle. Thus $(f(x),f(y),f(z))$ is not a $\Lambda$-triangle, and so either $f(z) > f(x)+f(y)$ or $f(z) < |f(x)-f(y)|$. Similarly for $z <|x-y|$. Now, to prove (b), assume toward a contradiction that there are $z_1, z_2\in \Delta$ with $z_1, z_2 >x+y$ such that $f(z_1)<|f(x)-f(y)|$ and $f(z_2)>f(x)+f(y)$. Without loss of generality assume $z_1<z_2$ (the argument for the case $z_1>z_2$ is similar). By the continuity of $f$ and the density of both $\Delta$ and $\Lambda$, we may find $z_1=a_1,\dots, a_n=z_2$ in $\Delta$ such that $a_k>x+y$ and $|f(a_{k+1})-f(a_k)| < f(x)+f(y)-|f(x)-f(y)|$ for all $1\leq k\leq n$. Then there must be an $a_k$ with $|f(x)-f(y)| \leq f(a_k) \leq f(x)+f(y)$, contradicting $a_k>x+y$. This proves (b). A similar argument for $z_1, z_2 < |x-y|$ proves (c). 

        Since $f$ is a bijection from $\Delta$ to $\Lambda$, the claimed inclusions are actually equalities. By considering sufficiently small $\delta_i$, we see that $f(\Delta_{<|x-y|}) = \Lambda_{<|f(x)-f(y)|}$, and thus $f(\Delta_{>x+y}) = \Lambda_{>f(x)+f(y)}$. By the continuty of $f$ again, we have $f(x+y) =f(x)+f(y)$.   
   
Finally, we argue that for any $x,t \in \Delta$, we have $f(x)/x=f(t)/t$. This proves (2). 


    For each $k > 0 $, find $0< \delta_k < 1/2^k$ by the continuity of $f$ such that for all $z\in \Delta$, $|t-z|<\delta_k$ implies $|f(t)-f(z)| < 1/k$. Now for each $k$, find positive rational number $n_k/m_k$ such that $0< t-(n_k/m_k)x < \delta_k/2$. Note that by the density of $\Delta$ and bijectivity of $f$, $\delta_k \to 0$ as $k \to \infty$. Thus $n_k/m_k \to t/x$ as $k \to \infty$. 

    Fix $k$. By the density of $\Delta$, find an increasing sequence $(t_j)_{j\in \mathbb{N}}$ with $t_j \in \Delta$ and $t_j \to x/m_k$ as $j \to\infty$. Then $m_kt_j \to x$ as $j \to \infty$, and $m_kt_j < x \in \Delta$. Thus $m_kt_j\in \Delta$ and $m_kf(t_j)=f(m_kt_j) \to f(x)$, which gives $f(t_j) \to f(x)/m_k$. On the other hand, $n_kt_j < (n_k/m_k)x < t \in \Delta$ and $n_kt_j \to (n_k/m_k)x$ as $j\to\infty$. Hence $n_kt_j\in \Delta$ and $f(n_kt_j)=n_kf(t_j) \to (n_k/m_k)f(x)$ as $j \to \infty$. 
    Thus we can find a $j$ such that $|n_kt_j-(n_k/m_k)x|<\delta_k/2$ and $|f(n_kt_j)-(n_k/m_k)f(x)|<1/k$. For this fixed $j$, we have $|t-n_kt_j|\leq |t-(n_k/m_k)x|+|n_kt_j-(n_k/m_k)x|<\delta_k$. Thus $|f(t)-f(n_kt_j)|<1/k$, and consequently $|f(t)-(n_k/m_k)f(x)|<2/k$. 

    Letting $k \to \infty$, we get $f(t)=(t/x)f(x)$ and thus $f(x)/x=f(t)/t$. 
\end{proof}

\begin{thebibliography}{99}
\bibitem{EGLMM}
M. Etedadialiabadi, S. Gao, F. Le Ma\^{\i}tre, and J. Melleray,
\textit{Dense locally finite subgroups of automorphism groups of ultraextensive spaces},
Adv. Math. {\bf 391} (2021), 107966.


\bibitem{GaoBook}
S. Gao,
Invariant Descriptive Set Theory,
Monographs and Textbooks in Pure and Applied Mathematics, vol. 293,
CRC Press, 2009.

\end{thebibliography}
\end{document}
